\begin{afterword}
\summaryhead

At thirteen or fourteen, the author thought Cantor's diagonal argument could be disproved. Every whole number can be written in binary, and the positions of its 1s form a set of whole numbers, so it seemed that the whole numbers could be matched with all such sets. A week later the author applied the diagonal argument to this matching and found a set it missed, the binary number $\ldots1111$.

The first impulse was to try again. Then the author saw that the failed attempt gave no more reason to doubt this theorem than any other, and that holding on was the road to becoming a crank who writes angry letters in green ink. The author let it go, and reflects that it was an unfair test for a child.

The lesson: the mistake was simply a mistake, not partly right and not a sign of virtue. Clinging to a thought you would never have had if you had been wiser is what makes a crackpot, even if no one ever hears of it. Say ``oops'' and move on.

\responsehead

The post tells a real mistake in enough detail to check, and the mathematics holds up. The binary pairing matches the whole numbers with the finite sets of whole numbers. The counterexample the author found, $\ldots1111$, is exactly the set the diagonal argument produces for this pairing. The turn of the story is also sound reasoning: a failed disproof leaves a proof where it was, so there was no more reason to doubt this theorem than any other.

The trouble is in how far the lesson reaches. A disproof is an unusually clean case. A proof with a gap proves nothing, so the insistence that the author was not ``the tiniest fraction right'' is exact here. Many mistaken beliefs are not like proofs: a forecast, a theory or a plan can be partly right. The post allows as much when it says that ``Not every cloud has a silver lining,'' which implies that some do. It does not say how to tell which kind of mistake one is facing.

The general rule has a related limit. It applies to ``a thought you would never have thought if you had been wiser,'' and a thought can be identified as one only after the mistake is found. The author's case was settled by a proof within a week. The situation in which a reader most needs advice, holding a view against the consensus with no quick way to check it, is the one the rule does not address. The post's subject is what to do after the mistake is known, and on that subject it is clear. It does not reach the earlier question.

In short: a candid account of a real mathematical mistake, with the mathematics right, whose rule for letting go works cleanly for a failed proof and gives no help until the mistake has been found.
\end{afterword}
